\chapter*{Tóm tắt}
%\markboth{TÓM TẮT}
Sự phát triển nhanh chóng của \ac{iot} đã mang lại nhiều tiện ích trong các lĩnh vực như nhà thông minh, công nghiệp, y tế và giao thông. Tuy nhiên, đi kèm là nguy cơ ngày càng gia tăng về các mối đe dọa an ninh mạng, đặc biệt khi thiết bị \ac{iot} thường hạn chế tài nguyên tính toán và bảo mật kém. Trong bối cảnh đó, \ac{ids} là một hướng phòng thủ nhằm phát hiện và ngăn chặn các hành vi tấn công bất thường vào hệ thống mạng \ac{iot}.

Luận văn này tập trung nghiên cứu, thiết kế và triển khai một hệ thống \ac{ids} dựa trên nền tảng xử lý dữ liệu lớn Apache Spark. Spark có khả năng xử lý dữ liệu phân tán, tốc độ cao và hỗ trợ hiệu quả các thuật toán học máy qua thư viện MLlib. Hệ thống phát hiện hành vi xâm nhập trong dữ liệu mạng \ac{iot} bằng các mô hình học máy truyền thống, học tổ hợp và học sâu MLP (cùng một autoencoder đóng vai trò cổng phát hiện bất thường ở tầng biên).

Luận văn phân tích đặc điểm dữ liệu mạng \ac{iot}, xử lý dữ liệu đầu vào từ tập dữ liệu công khai CICIDS2017, huấn luyện và đánh giá mô hình theo các chỉ số độ chính xác (Accuracy), precision, recall và F1-score, đồng thời ghi nhận thời gian huấn luyện và kích thước mô hình. Quá trình đánh giá áp dụng quy trình loại trừ rò rỉ dữ liệu (loại đặc trưng rò rỉ nhãn, loại bỏ trùng lặp ở mức đặc trưng) nhằm bảo đảm kết quả phản ánh đúng hiệu năng mô hình. Bên cạnh đó, hệ thống được kiểm thử hiệu năng trên cụm thiết bị biên nhằm khảo sát khả năng mở rộng và tính khả thi triển khai thực tiễn.

Dưới cấu hình đánh giá loại trừ rò rỉ dữ liệu, lựa chọn đặc trưng SHAP Top-30 giữ được hiệu năng gần tương đương cấu hình đầy đủ trong khi giảm khoảng 60\% số đặc trưng: Random Forest đạt F1 0{,}9955 và \ac{xgb} đạt 0{,}9970 trên CICIDS2017. Các khác biệt giữa phương pháp được kiểm chứng bằng kiểm định hoán vị ghép cặp và khoảng tin cậy hiệu chỉnh Nadeau--Bengio. Học tổ hợp (Hybrid Bagging 0{,}9966) không vượt được mô hình đơn tốt nhất, trong khi chi phí huấn luyện và dung lượng mô hình cao hơn nhiều lần. Trên cụm hai thiết bị biên Jetson Orin Nano Super, kiến trúc tách pipeline (cổng autoencoder trên một bo, bộ phân loại PySpark trên bo còn lại) duy trì được 22{,}4 luồng/giây, gấp khoảng 9 lần một nút (2{,}5 luồng/giây), với độ trễ đầu cuối p95 khoảng 7\,giây (đã gồm thời gian chờ hàng đợi) và năng lượng giảm từ 3{,}04 xuống 0{,}60\,J mỗi phán quyết. Cái giá là cổng bỏ qua khoảng 30\% luồng tấn công ở ngưỡng đã chọn. Cùng mô hình chạy qua ONNX Runtime nhanh hơn PySpark khoảng 7.200 lần, cho thấy nút thắt nằm ở engine suy luận chứ không ở mô hình. Khía cạnh xử lý phân tán được kiểm chứng ở mức minh chứng khái niệm (proof-of-concept), chưa khẳng định quy mô dữ liệu lớn thực thụ. Kết quả cho thấy một quy trình phát hiện xâm nhập thời gian thực khả thi trên thiết bị biên, đồng thời là cơ sở để kết hợp học máy với nền tảng dữ liệu lớn cho bảo mật \ac{iot}.